% ANEXO

\begin{anexosenv}
%\partanexos % Folha divisória (comentar caso não deseje que seja exibida)

% Primeiro anexo
\chapter{Nome do primeiro anexo} % edite para alterar o título deste anexo
\label{chap:anexoA}

Alguns exemplos de anexos: Mapas e documentos cartográficos; Leis, estatutos e regulamentos que esclareçam as condições jurídicas da pesquisa; Textos e reportagens na íntegra. Inserir aqui o arquivo do anexo que não precisa de texto introdutório, apenas o título na parte superior. 

Lembre-se que a diferença entre apêndice e anexo diz respeito à autoria do texto e/ou material ali colocado.

Caso o material ou texto suplementar ou complementar seja de sua autoria, então ele deverá ser colocado como um apêndice. Porém, caso a autoria seja de terceiros, então o material ou texto deverá ser colocado como anexo.

Caso seja conveniente, podem ser criados outros apêndices/anexos para o seu trabalho acadêmico. Basta recortar e colar este trecho neste mesmo documento. Lembre-se de alterar o "label".

% Segundo anexo
\chapter{Nome do segundo anexo}
\label{chap:anexoB}

Conteúdo do segundo anexo

\end{anexosenv}